\documentclass[11pt]{article}
\usepackage[margin=0.9in]{geometry}
\usepackage{amsmath,amssymb,booktabs,array,hyperref,fancyhdr}
\usepackage[T1]{fontenc}
\usepackage{lmodern}
\hypersetup{colorlinks=true,linkcolor=black,urlcolor=blue}
\setlength{\parindent}{0pt}
\setlength{\parskip}{6pt}
\allowdisplaybreaks
\pagestyle{fancy}\fancyhf{}\fancyhead[L]{Distribution-based pricing: mathematical specification}
\fancyfoot[C]{\thepage}\setlength{\headheight}{14pt}
\newcommand{\E}{\mathbb{E}}
\newcommand{\ind}{\mathbf{1}}
\newcommand{\dd}{\,\mathrm{d}}
\newcommand{\Cov}{\operatorname{Cov}}
\title{Distribution-Based Pricing\\\large Information, pricing paths, and welfare accounting}
\author{}\date{November 2026}
\begin{document}
\maketitle
\vspace{-2em}
This note records the mathematical objects behind \texttt{supply\_review.pdf}.
The price-setting equations include the corrected persistent-type introductory
objective. Terminal prices retain fixed beliefs: there is no learning.
Outcome accounting and the introductory objective use the same joint
type--history probabilities. No numerical
optima or interpretation of the quintile comparisons are supplied.

\section{Primitives and information}
There are two periods, $t=1,2$, one inside product $j$, and an outside option
$0$. Let $y_t\in\{j,0\}$ denote the choice in period $t$ and let
$h_0\in\{j,0\}$ denote the given pre-period-1 choice. The state entering
period 2 is $h_1=y_1$. A consumer's type $\gamma$ is persistent.
Parameters satisfy $\alpha<0$, $\delta\in[0,1]$; $\beta$ is the inside
intercept. Marginal cost is zero in the current packet. Habit stocks and
other estimated heterogeneity are not advanced in this two-period exercise.

Let $F$ be the known population distribution of $\gamma$ and define
\begin{align}
 b_k&=F^{-1}(k/5),\qquad k=0,\ldots,5,\\
 B_k&=(b_{k-1},b_k],\qquad
 \pi_k=\Pr(\gamma\in B_k)=\tfrac15,\\
 F_k(A)&=\Pr(\gamma\in A\mid\gamma\in B_k).
\end{align}
Endpoints have zero probability for the continuous distributions here.
The uniform firm knows $F$, every $F_k$, and every $\pi_k$; it cannot place
the individual in a quintile. The informed firm additionally observes the
individual's quintile $k$, but not their exact $\gamma$.
Thus, for any integrable $g$,
\begin{equation}
 \boxed{\E_F[g(\gamma)]
 =\sum_{k=1}^5\pi_k\E_{F_k}[g(\gamma)].}
 \label{eq:mixture}
\end{equation}
Write $G=F$ for uniform pricing and $G=F_k$ for quintile-informed pricing.
Uniform means common across unobserved quintiles at a given observed
history; it does not mean a price constant across time or histories.

\section{Utility and purchase probabilities}
The deterministic utility levels used in welfare accounting are
\begin{align}
 V_j(p,\gamma,h)&=\beta+\alpha p+\gamma\ind\{h=0\},\\
 V_0(\gamma,h)&=\gamma\ind\{h=j\}.
\end{align}
Add independent type-I extreme-value choice shocks with unit scale.
Consumers choose myopically. Define $a_j=-1$, $a_0=1$, and
\begin{align}
 z_h(p,\gamma)&=V_j-V_0=\beta+\alpha p+a_h\gamma,\\
 \Lambda(z)&=\frac{1}{1+e^{-z}},\\
 q_h(p,\gamma)&=\Pr(y_t=j\mid p,\gamma,h)
 =\Lambda(z_h(p,\gamma)).
\end{align}
The type-level price derivatives are
\begin{align}
 q_h'(p,\gamma)&=\alpha q_h(p,\gamma)[1-q_h(p,\gamma)],\\
 q_h''(p,\gamma)&=\alpha^2q_h(p,\gamma)[1-q_h(p,\gamma)]
 [1-2q_h(p,\gamma)].
\end{align}
For a belief distribution $G$, aggregate demand and its derivatives are
\begin{align}
 S_h^G(p)&=\int q_h(p,\gamma)\dd G(\gamma),\\
 (S_h^G)'(p)&=\alpha\int q_h(p,\gamma)[1-q_h(p,\gamma)]\dd G(\gamma),\\
 (S_h^G)''(p)&=\alpha^2\int q_h(p,\gamma)[1-q_h(p,\gamma)]
 [1-2q_h(p,\gamma)]\dd G(\gamma).
\end{align}
In particular,
\begin{align}
 S_h^F(p)&=\sum_k\pi_k S_h^{F_k}(p),&
 (S_h^F)'(p)&=\sum_k\pi_k(S_h^{F_k})'(p).
\end{align}
These are integrals of demand, not demand evaluated at mean gamma:
\begin{equation}
 S_h^G(p)=\E_G[\Lambda(\beta+\alpha p+a_h\gamma)],
 \qquad
 \widetilde S_h^G(p)=\Lambda(\beta+\alpha p+a_h\E_G[\gamma]).
\end{equation}
The code uses $S$, not $\widetilde S$.

\section{Second-period pricing under the implemented beliefs}
The current solver holds $G$ fixed after either purchase history. For
$h\in\{j,0\}$ it solves
\begin{align}
 R_h^G(p)&=pS_h^G(p),\\
 p_{2h}^G&\in\arg\max_{p\geq0}R_h^G(p),\\
 W_h^G&=R_h^G(p_{2h}^G).
\end{align}
An interior solution satisfies
\begin{align}
 0&=S_h^G(p_{2h}^G)+p_{2h}^G(S_h^G)'(p_{2h}^G),\label{eq:terminalfoc}\\
 p_{2h}^G
 &=\frac{\E_G[q_h]}{(-\alpha)\E_G[q_h(1-q_h)]},\label{eq:markup}
\end{align}
where the probabilities in \eqref{eq:markup} are evaluated at the optimum.
Equivalently, the absolute price elasticity
$-p(S_h^G)'(p)/S_h^G(p)$ equals one there. A local maximum also requires
\begin{equation}
 2(S_h^G)'(p)+p(S_h^G)''(p)\leq0.
\end{equation}
For mixtures, stationary points must be compared with one another and the
search boundaries; the first-order condition alone does not select a global optimum.

For any candidate price, define the quintile contribution
\begin{equation}
 D_{kh}(p)=S_h^{F_k}(p)+p(S_h^{F_k})'(p).
\end{equation}
The uniform revenue derivative is exactly
\begin{equation}
 (R_h^F)'(p)=\sum_{k=1}^5\pi_k D_{kh}(p).\label{eq:contrib}
\end{equation}
At a Q3 interior optimum $\widehat p=p_{2h}^{F_3}$, the condition for the
same price also to be a uniform stationary point is
\begin{equation}
 \sum_{k\ne3}\pi_k D_{kh}(\widehat p)=0.
\end{equation}
This is the marginal-revenue comparison relevant to the proposed proximity
of the uniform and 40--60\% prices. Equality of mean gamma is a separate condition.

Another exact expression for the markup denominator is
\begin{equation}
 \E_G[q_h(1-q_h)]
 =S_h^G(p)[1-S_h^G(p)]-\operatorname{Var}_G(q_h(p,\gamma)).
\end{equation}
All objects here depend on the candidate price.

\section{First-period pricing in the current solver}
Given the no-learning terminal prices $p_{2j}^G,p_{20}^G$, define
\begin{align}
 r_j^G(\gamma)&=p_{2j}^Gq_j(p_{2j}^G,\gamma),\\
 r_0^G(\gamma)&=p_{20}^Gq_0(p_{20}^G,\gamma),\\
 \Delta r^G(\gamma)&=r_j^G(\gamma)-r_0^G(\gamma).
\end{align}
The corrected first-period objective is
\begin{align}
 J_{h_0}^G(p)&=\E_G\!\left[pq_{h_0}(p,\gamma)+\delta\{
 q_{h_0}(p,\gamma)r_j^G(\gamma)
 +[1-q_{h_0}(p,\gamma)]r_0^G(\gamma)\}\right],\\
 p_1^G(h_0)&\in\arg\max_p J_{h_0}^G(p).
\end{align}
The first-period search permits a negative introductory price. Terminal prices
are fixed while choosing $p$ and do not use posterior type weights. Thus
\begin{align}
 J_{h_0}^G(p)&=\delta\E_G[r_0^G]
 +\E_G[q_{h_0}(p,\gamma)(p+\delta\Delta r^G(\gamma))],\\
 (J_{h_0}^G)'(p)&=\E_G[q_{h_0}+q_{h_0}'(p+\delta\Delta r^G)],\\
 (J_{h_0}^G)''(p)&=\E_G[2q_{h_0}'+q_{h_0}''(p+\delta\Delta r^G)].
\end{align}
Writing $q=q_{h_0}(p_1^G,\gamma)$, an interior optimum satisfies
\begin{equation}
 \boxed{p_1^G(h_0)=\frac{\E_G[q]}{(-\alpha)\E_G[q(1-q)]}
 -\delta\frac{\E_G[q(1-q)\Delta r^G]}{\E_G[q(1-q)]}.}
\end{equation}
The continuation adjustment is type-specific and enters with weights
proportional to $q(1-q)$. It is not replaced by the unweighted mean
$\E_G[\Delta r^G]$. The price and probabilities on the right depend on one
another; this is an implicit condition. Global candidate checks remain necessary.

The uniform derivative decomposes by quintile as
\begin{align}
 (J_{h_0}^F)'(p)
 &=\sum_k\pi_k\E_{F_k}[q_{h_0}+q_{h_0}'(p+\delta\Delta r^F)].
\end{align}
At an interior Q3 optimum $\widehat p=p_1^{F_3}(h_0)$, its Q3 summand equals
\begin{equation}
 \pi_3\delta\E_{F_3}[q_{h_0}'(\widehat p,\gamma)
 (\Delta r^F(\gamma)-\Delta r^{F_3}(\gamma))].
\end{equation}
Uniform and informed continuation policies differ, even when evaluated on
the same quintile. This decomposition uses each policy's actual terminal offers.

\section{Normal distributions and numerical integration}
In the packet $F=N(\mu,\sigma^2)$. The quintile boundaries are
$b_k=\mu+\sigma\Phi^{-1}(k/5)$.
For $0<k<5$, $\Phi$ is the standard normal CDF and $\phi$ its density;
$b_0=-\infty$, $b_5=+\infty$. Conditional means are
\begin{equation}
 \E[\gamma\mid B_k]=\mu+5\sigma\left[
 \phi(\Phi^{-1}((k-1)/5))-\phi(\Phi^{-1}(k/5))\right].
\end{equation}
This formula describes a mean; the solver retains the within-bin distribution.
For integration, let $(x_\ell,w_\ell)_{\ell=1}^L$ be Gauss--Legendre nodes
and weights on $[-1,1]$, with $L=21$ and $\sum_\ell w_\ell=2$. Define
\begin{align}
 u_{k\ell}&=\frac{k-1+(x_\ell+1)/2}{5},\\
 \gamma_{k\ell}&=\mu+\sigma\Phi^{-1}(u_{k\ell}),\\
 \omega_{k\ell}&=\frac{w_\ell}{10},&
 \widetilde\omega_{k\ell}&=5\omega_{k\ell}=\frac{w_\ell}{2}.
\end{align}
Thus
\begin{align}
 \E_F[g]&\approx\sum_{k,\ell}\omega_{k\ell}g(\gamma_{k\ell}),\\
 \E_{F_k}[g]&\approx\sum_\ell\widetilde\omega_{k\ell}g(\gamma_{k\ell}),\\
 \sum_\ell\omega_{k\ell}&=\tfrac15,&
 \sum_\ell\widetilde\omega_{k\ell}&=1.
\end{align}
Each candidate price is evaluated at every integration node before summing.
The uniform solver uses all five bins with their known weights.

\section{Histories and the information set}
For the empirical analogue, let $t_i^0$ denote household $i$'s first actual
observed yogurt basket. Zero-yogurt weeks do not define the initial basket.
For focal flavor $f$, its observed initial state is
\begin{equation}
 h_{i0}^f=\begin{cases}j,&\text{positive focal-flavor purchases at }t_i^0,\\
 0,&\text{no focal-flavor purchases at }t_i^0.\end{cases}
\end{equation}
Index the first observed basket by purchase occasion 0. The counterfactual
pricing periods are purchase occasions 1 and 2: the next two consecutive
purchase weeks, which need not be consecutive calendar weeks. Calendar gaps
and zero-yogurt weeks do not advance this time index. Discounting uses $\delta$
per pricing period, not per elapsed calendar week.
Their offers are optimized and their purchases are endogenous. Later observed
choices are not imposed as counterfactual outcomes. An initial basket with
unknown focal status is excluded and audited; a later usable occasion does
not replace the first occasion.

The maintained type belief is the same population $G$ under both initial
histories. Initial history affects utility, not the distribution of types
used to choose the price. Write $N_h^f=\sum_i\ind\{h_{i0}^f=h\}$ and
$N^f=N_j^f+N_0^f$. For any outcome $M$ and policy $r$, the observed-initialized
population average is
\begin{equation}
 \overline M_f^r=\sum_{h\in\{j,0\}}\frac{N_h^f}{N^f}
 \E_F[M_f^r(\gamma,h)].
\end{equation}
The informed expression uses the appropriate quintile-specific policy inside
the integral. Observed initial frequencies weight the combined population;
they do not change the optimum conditional on a given $h_0$ when all primitives
and type beliefs conditional on that state are held fixed. Quintiles are
population probability bins, not recovered household posterior ranks.

The local population-input figures evaluate the two conditional states but do
not read $t_i^0$ or $N_h^f$. The observed-initialization version is generated
on the research server by \texttt{run\_observed\_state\_theory.py}.

For a fixed first-period policy, write
$q_1(\gamma)=q_{h_0}(p_1,\gamma)$ and $s_1=\E_G[q_1]$.
The actual conditional composition of the two outcome branches is
\begin{align}
 \dd G_j(\gamma;p_1,h_0)&=\frac{q_1(\gamma)\dd G(\gamma)}{s_1},\\
 \dd G_0(\gamma;p_1,h_0)&=\frac{[1-q_1(\gamma)]\dd G(\gamma)}{1-s_1}.
\end{align}
For the uniform population the actual branch-specific quintile shares are
\begin{align}
 \Pr(B_k\mid y_1=j,h_0,p_1)
 &=\frac{\pi_k\E_{F_k}[q_1]}{\sum_m\pi_m\E_{F_m}[q_1]},\\
 \Pr(B_k\mid y_1=0,h_0,p_1)
 &=\frac{\pi_k\E_{F_k}[1-q_1]}{\sum_m\pi_m\E_{F_m}[1-q_1]}.
\end{align}
These conditional compositions are accounting identities, not beliefs used
to choose the terminal offers. The maintained no-learning restriction is
\begin{equation}
 G_{\mathrm{price},j}=G_{\mathrm{price},0}=G.
\end{equation}
The introductory objective anticipates persistent-type outcomes under that
fixed terminal rule. Terminal prices do not use the conditional compositions
above, depend on the initial offer, or acquire additional information about
quintile membership. This is the specified fixed-belief policy restriction,
not Bayesian backward induction.

\section{Persistent-type outcome accounting}
For a given price triple $(p_1,p_{2j},p_{20})$, define
\begin{align}
 q_1(\gamma)&=q_{h_0}(p_1,\gamma),&
 q_{2j}(\gamma)&=q_j(p_{2j},\gamma),&
 q_{20}(\gamma)&=q_0(p_{20},\gamma),\\
 q_2(\gamma)&=q_1q_{2j}+(1-q_1)q_{20},\\
 r_j(\gamma)&=p_{2j}q_{2j}(\gamma),&
 r_0(\gamma)&=p_{20}q_{20}(\gamma).
\end{align}
Actual expected revenue for persistent types is
\begin{align}
 R_{\mathrm{act}}^G
 &=\E_G\!\left[p_1q_1+\delta\{q_1r_j+(1-q_1)r_0\}\right].\label{eq:actual}
\end{align}
The corrected introductory objective is exactly
\begin{equation}
 J_{h_0}^G(p_1)=R_{\mathrm{act}}^G.
\end{equation}
The implementation verifies this equality for each solved policy. Terminal
prices still maximize static revenue using the unchanged distribution $G$;
the correction does not introduce learning or jointly reoptimize terminal
prices for the selected consumers in each branch.

For reference, the superseded product-of-means objective was
\begin{equation}
 J_{\mathrm{old}}^G=p_1\E_G[q_1]+\delta\{\E_G[q_1]\E_G[r_j]
 +[1-\E_G[q_1]]\E_G[r_0]\}.
\end{equation}
Consequently,
\begin{equation}
 \boxed{R_{\mathrm{act}}^G-J_{\mathrm{old}}^G
 =\delta\Cov_G(q_1,r_j-r_0).}\label{eq:cov}
\end{equation}
Equation \eqref{eq:cov} is an accounting identity for any such fixed policy.
The current solver integrates products before averaging; it no longer uses
$J_{\mathrm{old}}^G$ to select the introductory price.

Even retaining the fixed terminal prices, the derivative of actual revenue
with respect to the introductory price is
\begin{equation}
 \frac{\partial R_{\mathrm{act}}^G}{\partial p_1}
 =\E_G[q_1]+p_1\E_G[q_1']
 +\delta\E_G[q_1'(r_j-r_0)].
\end{equation}
This is a derivative holding the no-learning terminal policy fixed.
Conditional terminal revenue for the buyer
branch can equivalently be written
\begin{equation}
 \E_G[q_1r_j]=s_1\E_{G_j}[r_j],
\end{equation}
with the analogous nonbuyer expression.

\section{What the displayed price paths average}
The main no-habit results packet now displays four decision panels:
$p_1^r(j)$, $p_1^r(0)$, $p_{2j}^r$, and $p_{20}^r$, each against quintile.
These are optimized offers without history averaging. The following
definitions describe the supplemental expected-price figures and saved tables.
For policy $r\in\{U,I_k\}$ and initial history $h_0$, the individual-type
expected second-period offer is
\begin{equation}
 e_2^r(\gamma\mid h_0)=q_{h_0}(p_1^r,\gamma)p_{2j}^r
 +[1-q_{h_0}(p_1^r,\gamma)]p_{20}^r.
\end{equation}
This averages offers over purchase histories, not transaction prices
conditional on purchasing in period 2. It is not discounted.
The informed quintile path is
\begin{equation}
 \left(p_1^{I_k}(h_0),\;\E_{F_k}[e_2^{I_k}(\gamma\mid h_0)]\right).
\end{equation}
The single uniform line displayed in each panel is
\begin{equation}
 \left(p_1^U(h_0),\;\E_F[e_2^U(\gamma\mid h_0)]\right).
\end{equation}
The uniform-policy expectation specifically for quintile $k$ is instead
\begin{equation}
 \E_{F_k}[e_2^U]=\E_{F_k}[q_1^U]p_{2j}^U
 +[1-\E_{F_k}[q_1^U]]p_{20}^U,
\end{equation}
and the plotted population expectation satisfies
\begin{equation}
 \E_F[e_2^U]=\sum_k\pi_k\E_{F_k}[e_2^U].
\end{equation}
Welfare comparisons use the same quintile under both policies, even though
the price-path panel displays one population-averaged uniform line.

\section{Consumer surplus and total welfare}
The expected maximum utility, in dollars and omitting the common shock
constant, is
\begin{align}
 L_h(p,\gamma)
 &=\frac{\log(e^{V_j}+e^{V_0})}{-\alpha}\\
 &=\frac{\log(1+e^{z_h(p,\gamma)})+\gamma\ind\{h=j\}}{-\alpha}.
\end{align}
Thus the period-specific and total type-level quantities are
\begin{align}
 CS_1^r(\gamma,h_0)&=L_{h_0}(p_1^r,\gamma),\\
 CS_2^r(\gamma,h_0)&=q_1^r(\gamma)L_j(p_{2j}^r,\gamma)
 +[1-q_1^r(\gamma)]L_0(p_{20}^r,\gamma),\\
 CS^r(\gamma,h_0)&=CS_1^r+\delta CS_2^r,\\
 R^r(\gamma,h_0)&=p_1^rq_1^r+\delta\{q_1^rp_{2j}^rq_{2j}^r
 +(1-q_1^r)p_{20}^rq_{20}^r\},\\
 TS^r(\gamma,h_0)&=CS^r(\gamma,h_0)+R^r(\gamma,h_0).
\end{align}
These use the same persistent $\gamma$ inside every history probability and
utility term. Surplus is experienced utility under myopic choices, not the
value function of a forward-looking consumer.

For comparison, deleting the outside utility level from both alternatives
leaves choice probabilities unchanged and gives
\begin{align}
 \widetilde L_h(p,\gamma)&=\frac{\log(1+e^{z_h(p,\gamma)})}{-\alpha},\\
 CS^r-\widetilde{CS}^r
 &=\frac{\gamma}{-\alpha}\{\ind\{h_0=j\}+\delta q_1^r(\gamma)\}.
\end{align}
For the same initial history, the difference between policy comparisons is
\begin{equation}
 (CS^I-CS^U)-(\widetilde{CS}^I-\widetilde{CS}^U)
 =\frac{\delta\gamma}{-\alpha}(q_1^I-q_1^U).
\end{equation}
The packet retains $L_h$, including the outside utility level.

For $M\in\{CS,R,TS\}$, its conditional quintile mean and plotted change are
\begin{align}
 \overline M_k^r(h_0)&=\int M^r(\gamma,h_0)\dd F_k(\gamma),\\
 \Delta\overline M_k(h_0)&=\overline M_k^{I_k}(h_0)-\overline M_k^U(h_0),\\
 P_{M,k}(h_0)&=100\frac{\Delta\overline M_k(h_0)}{|\overline M_k^U(h_0)|}.
\end{align}
The denominator is the absolute baseline of the same metric, quintile, and
initial history. A zero baseline makes the percentage undefined. This is
a ratio of conditional means, not a mean of type-level percentage changes.

\section{Parameter inputs and implementation references}
The following are inputs, not optimized prices. All exercises use
$\delta=1$, zero marginal cost, five equally probable quintiles, and
21 probability-space integration nodes per quintile.
\begin{center}\small
\begin{tabular}{lrrrr}\toprule
Exercise & $\mu$ & $\sigma$ & $\alpha$ & $\beta$\\\midrule
Pure theory & 0 & 1.25 & $-1$ & 2\\
Earlier plain & $-0.357612$ & 1.130644 & $-0.802523$ & 0.698627\\
Earlier berry & $-0.357612$ & 1.130644 & $-0.802523$ & 0.402734\\
Earlier other & $-0.357612$ & 1.130644 & $-0.802523$ & 0.692726\\
Observed-week plain & $-0.058010$ & 1.460255 & $-1.838628$ & $-2.523088$\\
Observed-week berry & $-0.058010$ & 1.460255 & $-1.838628$ & $-2.498755$\\
Observed-week other & $-0.058010$ & 1.460255 & $-1.838628$ & $-1.730666$\\
No-habit plain & $-0.838193$ & 2.099165 & $-3.530427$ & $-1.285109$\\
No-habit berry & $-0.838193$ & 2.099165 & $-3.530427$ & $-1.619406$\\
No-habit other & $-0.838193$ & 2.099165 & $-3.530427$ & 0.016210\\
\bottomrule\end{tabular}
\end{center}
Full-precision inputs are saved in each exercise's \texttt{settings.json}.
The current empirical analogue uses the no-habit rows. Its source is
\par{\small\path{separate_experiments_50281568/novelty_only/point_estimates.csv}.}\par
For the earlier-fit exercises the packet imports saved theoretical-node
outcomes. Pure theory and observed-week exercises are solved by
\texttt{quintile\_theory.run}. No empirical records are needed for this note.

\begin{description}
\item[Pricing solver:]
\path{Code/final_simple/supply/quintile_theory.py},\\
functions \texttt{run}
and \texttt{solve\_distribution}.
\item[Earlier-fit pricing:]
\path{Code/final_simple/demand/quantified_theory.py},\\
\texttt{gamma\_bin\_distribution} and \texttt{compare\_bins};\\
\path{Code/final_simple/demand/gamma_calibration.py}, \texttt{uncertainty}.
\item[Outcome accounting:]
\path{Code/final_simple/demand/binary_welfare.py}, \texttt{evaluate\_policy}.
\item[Displayed averaging:]
\path{Code/final_simple/demand/plot_supply_bins.py}, \texttt{plot\_bins}.
\item[Packet assembly:]
\path{Code/final_simple/supply/build_supply_review.py}, \texttt{build}.
\end{description}
\end{document}
